\documentclass[12pt]{article}

\usepackage[utf8]{inputenc}
\usepackage[romanian]{babel}
\usepackage{amsmath, amssymb, amsthm}
\usepackage{geometry}
\geometry{a4paper, margin=2.5cm}

\title{Exemple și Contraexemple}
\author{}
\date{}

\theoremstyle{definition}
\newtheorem{definition}{Definiție}

\begin{document}

\maketitle

\begin{definition}[Exemple și contraexemple]
Exemplele sunt cazuri particulare cu funcții specifice în activitatea matematică și didactică. Ele:
\begin{enumerate}
    \item probează enunțuri existențial–afirmative;
    \item arată consistența teoriilor;
    \item contribuie la o mai bună înțelegere a noilor noțiuni;
    \item facilitează sesizarea diferențelor dintre concepte apropiate.
\end{enumerate}
\end{definition}

\end{document}
